\chapter{Vector Spaces}
\section{Definitions}
\begin{definitionenv}
		A \textbf{field} $\mathcal{K}$ is a set equipped with addition and multiplication, denoted by $(\mathcal{K},+,\cdot)$, satisfying the following axioms:
		\begin{enumerate}
			\item \textbf{Associative Laws}: let $a,b,c \in \mathcal{K}$, $$a+b+c=a+(b+c), \qquad a\cdot b\cdot c=a\cdot (b\cdot c).$$
			\item \textbf{Existence of Identity Elements}: let $a\in \mathcal K$, $$a+0_{\mathcal{K}}=a,\qquad a\cdot 1_{\mathcal{K}}=a,$$ where $0_{\mathcal{K}}$ and $1_\mathcal{K}$ is called the additive identity and multiplicative identity.
			\item \textbf{Existence of Negative}: for any $a\in \mathcal{K}$, there exists one element in $\mathcal{K}$ denoted by $(-a)$ such that $$a+(-a)=0_\mathcal{K}.$$
			\item \textbf{Existence of Reciprocal}: for any $a\in \mathcal{K}$ and $a\neq 0_\mathcal{K}$, there exists one element in $\mathcal{K}$ denoted by $a^{-1}$ such that $$a\cdot a^{-1} = 1_\mathcal{K}.$$ 
			\item \textbf{Commutative Laws}: let $ a,b\in \mathcal{K}$, $$a+b=b+a, \qquad a\cdot b=b\cdot a.$$
			\item \textbf{Distributive Law}: let $a,b,c \in \mathcal{K}$, $$a\cdot (b+c)=a\cdot b + a\cdot c.$$
		\end{enumerate}
	\end{definitionenv}	
	\begin{remarkenv}
		Closure is contained in the defined addition and multiplication.
	\end{remarkenv}
    \begin{notationenv}
        For $a,b \in \mathcal{K}$, we usually write $a-b$ to replace $a+(-b)$ and write $\dfrac{a}{b}$ or $a/b$ to replace $a\cdot b^{-1}$ when $b\neq 0_\mathcal{K}$.
    \end{notationenv}
	\begin{theoremenv}[Cancellation Law]
		Let $a,b,c \in \mathcal{K}$, if $a+c=b+c$, then $a=b$.
	\end{theoremenv}
	\begin{corollaryenv}
		$0_\mathcal{K}$ and $1_\mathcal{K}$ in $\mathbf{Field\ Axiom\ 2}$ is unique.
	\end{corollaryenv}
	\begin{exampleenv}
		$\mathbb{C,R,Q}$ are fields, but $\mathbb{N}$ is not a field because $2 \in \mathbb{N}$, but its reciprocal $\dfrac{1}{2} \notin \mathbb{N}$. This is the main difference between $\mathbb{Q}$ and $\mathbb{N}$.
	\end{exampleenv}
	\begin{definitionenv}\label{def_of_ss}
		Let $\mathcal{F}$ be a field. A subset $\mathcal{K} \subseteq \mathcal{F}$ is called a \textbf{subfield} of $\mathcal{F}$ if $\mathcal{K}$ is itself a field under the same operations of addition and multiplication as $\mathcal{F}$. Equivalently, $\mathcal{K}$ is a subfield of $\mathcal{F}$ if and only if the following conditions hold:
		\begin{enumerate}
			\item \textbf{Closure under Addition and Multiplication:} If $x,y$ are elements of $\mathcal{K}$, then  $x+y\in \mathcal{K}$ and $x\cdot y$ $\in \mathcal{K}$.
			\item \textbf{Existence of Additive and Multiplicative Inverses:} If $x\in \mathcal{K}$, then $(-x)\in \mathcal{K}$. If furthermore $x\neq 0$ , $x^{-1}$ $\in \mathcal{K}$.
			\item \textbf{Existence of Identity Elements:} $0_\mathcal{K}$ and $1_\mathcal{K}$ are in $\mathcal{K}$.
		\end{enumerate}
	\end{definitionenv}
	\begin{remarkenv}
		Definition~\ref{def_of_ss} provides a simpler way to prove a set to be a field. The third rule ensures that $\varnothing$ is not a field.
	\end{remarkenv}
	\begin{propositionenv}
		Let $\mathcal{K}$ be a field and $A,B$ are subfields of $\mathcal{K}$, then $A \cap B$ is also a subfield of $\mathcal{K}$.
	\end{propositionenv}
	\begin{propositionenv}
		Let $\mathcal{K}$ be a field and $A,B$ are subfields of $\mathcal{K}$, then $A \cup B$ is also a subfield of $\mathcal{K}$ if and only if $A \subseteq B$ or $B \subseteq A$.
	\end{propositionenv}
\begin{remarkenv}
	For any elements in a field, we call them \textbf{number}s or \textbf{scalar}s.
\end{remarkenv}
\begin{definitionenv}
	A \textbf{vector space} $V$ over a field $\mathcal{K}$ is a set equipped with two operations:
	vector addition and scalar multiplication, denoted by $(V, +, \cdot)$, satisfying the following axioms:

	\begin{enumerate}
		\item \textbf{Associative Laws:} for any $\mathbf{u}, \mathbf{v}, \mathbf{w} \in V$,
		\[
		\mathbf{u} + (\mathbf{v} + \mathbf{w}) = (\mathbf{u} + \mathbf{v}) + \mathbf{w}.
		\]

		\item \textbf{Existence of Identity Elements:}
		there exists a zero vector $\mathbf{0} \in V$ such that for any $\mathbf{v} \in V$,
		\[
		\mathbf{v} + \mathbf{0} = \mathbf{v}.
		\]

		\item \textbf{Existence of Additive Inverses:}
		for any $\mathbf{v} \in V$, there exists an element $-\mathbf{v} \in V$ such that
		\[
		\mathbf{v} + (-\mathbf{v}) = \mathbf{0}.
		\]

		\item \textbf{Commutative Law:}
		for any $\mathbf{u}, \mathbf{v} \in V$,
		\[
		\mathbf{u} + \mathbf{v} = \mathbf{v} + \mathbf{u}.
		\]

		\item \textbf{Distributive Laws:}
		for all $a, b \in \mathcal{K}$ and $\mathbf{u}, \mathbf{v} \in V$,
		\[
		a \cdot (\mathbf{u} + \mathbf{v}) = a \cdot \mathbf{u} + a \cdot \mathbf{v},
		\quad (a + b) \cdot \mathbf{v} = a \cdot \mathbf{v} + b \cdot \mathbf{v}.
		\]

		\item \textbf{Compatibility of Scalar Multiplication:}
		for all $a, b \in \mathcal{K}$ and $\mathbf{v} \in V$,
		\[
		(a b) \cdot \mathbf{v} = a \cdot (b \cdot \mathbf{v}).
		\]

		\item \textbf{Identity of Scalar Multiplication:}
		for all $\mathbf{v} \in V$,
		\[
		1_{\mathcal{K}} \cdot \mathbf{v} = \mathbf{v}.
		\]
	\end{enumerate}
\end{definitionenv}
\begin{theoremenv}
	The zero vector $\mathbf{0}$ is unique.
\end{theoremenv}
\begin{definitionenv}
	Let $V$ be a vector space over a field $\mathcal{K}$.
	A non-empty subset $W \subseteq V$ is called a \textbf{subspace} of $V$ if $W$ is itself a vector space over $\mathcal{K}$ under the same operations of addition and scalar multiplication as $V$.

	Equivalently, $W$ is a subspace of $V$ if and only if the following conditions hold for all $\mathbf{u}, \mathbf{v} \in W$ and $a \in \mathcal{K}$:
	\begin{enumerate}
		\item \textbf{Closure under Addition:}
		\[
		\mathbf{u} + \mathbf{v} \in W.
		\]
		\item \textbf{Closure under Scalar Multiplication:}
		\[
		a \cdot \mathbf{u} \in W.
		\]
	\end{enumerate}

\end{definitionenv}
\begin{notationenv}
	If $W$ is a subspace of $V$, we write $W \leq V$.
\end{notationenv}
\begin{remarkenv}
	We state $W$ to be a non-empty set at first so we needn't to say $\mathbf{0}\in W$.
\end{remarkenv}
\begin{remarkenv}
	For any elements in a vector space, we call them \textbf{vector}s.
\end{remarkenv}
\begin{exampleenv}
	If $\mathcal{F}$ is a field, then $\mathcal{F}^n$ is always a vector space over $\mathcal{F}$. Furthermore, $\mathcal{F}^n$ is a vector space over $\mathcal{K}$ whenever $\mathcal{K}\subseteq \mathcal{F}$ is a subfield of $\mathcal{F}$.
\end{exampleenv}
\begin{definitionenv}
	Let $V$ be a vector space over a field $\mathcal{K}$, and let vectors
	$\mathbf{v}_1, \mathbf{v}_2, \dots, \mathbf{v}_n \in V$.
	Any vector of the form
	\[
	\mathbf{v} = a_1 \mathbf{v}_1 + a_2 \mathbf{v}_2 + \cdots + a_n \mathbf{v}_n,
	\]
	where $a_1, a_2, \dots, a_n \in \mathcal{K}$, is called a \textbf{linear combination} of
	$\mathbf{v}_1, \mathbf{v}_2, \dots, \mathbf{v}_n$.
\end{definitionenv}
\begin{definitionenv}
	Let $V$ be a vector space over a field $\mathcal{K}$, and let
	$S = \{\mathbf{v}_1, \mathbf{v}_2, \dots, \mathbf{v}_n\} \subseteq V$.
	The set of all linear combinations of vectors in $S$ is called the
	\textbf{span} of $S$, denoted by
	\[
	\mathrm{span}(S) =
	\left\{
	a_1 \mathbf{v}_1 + a_2 \mathbf{v}_2 + \cdots + a_n \mathbf{v}_n
	\ \middle| \ a_1, a_2, \dots, a_n \in \mathcal{K}
	\right\}.
	\]
\end{definitionenv}
\begin{theoremenv}
	The span of $S$ is the smallest subspace of $V$ containing $S$. Particularly, span$(\varnothing)=\{\mathbf{0}\}=\text{span}\{\mathbf{0}\}$.
\end{theoremenv}
\begin{definitionenv}
	Let $V$ be a vector space over a field $\mathcal{K}$.
	A subset $S \subseteq V$ is called a \textbf{spanning set} of $V$ or a set of \textbf{generators}
	if $\mathrm{span}(S) = V$, that is, every vector in $V$ can be expressed
	as a linear combination of finitely many vectors from $S$.
\end{definitionenv}
\begin{definitionenv}
	Let $V$ be a vector space over a field $\mathcal{K}$, and let vectors
	$\mathbf{v}_1, \mathbf{v}_2, \dots, \mathbf{v}_n \in V$.
	The set $\{\mathbf{v}_1, \mathbf{v}_2, \dots, \mathbf{v}_n\}$ is said to be
	\textbf{linearly dependent} if there exist scalars
	$a_1, a_2, \dots, a_n \in \mathcal{K}$, not all zero, such that
	\[
	a_1 \mathbf{v}_1 + a_2 \mathbf{v}_2 + \cdots + a_n \mathbf{v}_n = \mathbf{0}.
	\]
	Otherwise, if this equation holds only when
	$a_1 = a_2 = \cdots = a_n = 0$, the set is said to be
	\textbf{linearly independent}.
\end{definitionenv}
\begin{remarkenv}
	$\{\mathbf{0}\}$ is linearly dependent, but $\varnothing$ is linearly independent.
\end{remarkenv}


\section{Bases}
\begin{definitionenv}
	Let $V$ be a vector space over a field $\mathcal{K}$.
	A subset $\mathcal{B} = \{\mathbf{v}_1, \mathbf{v}_2, \dots, \mathbf{v}_n\} \subseteq V$ is called a \textbf{basis} of $V$ if it satisfies the following two conditions:
	\begin{enumerate}
		\item $\mathcal{B}$ is linearly independent.
		\item $\mathcal{B}$ spans $V$, i.e. $\mathrm{span}(\mathcal{B}) = V$.
	\end{enumerate}
\end{definitionenv}
\begin{theoremenv}
	 Every vector in $V$ can be written uniquely as a linear combination of the vectors in $\mathcal{B}$.
\end{theoremenv}

\begin{remarkenv}
	Though we use set notation to denote a basis, the order $\emph{does}$ matter. For instance, two bases
	\[
	\mathcal{B}_1=\{(1,0),(0,1)\},\quad \mathcal{B}_2=\{(0,1),(1,0)\}
	\]
	are considered to be different. You will know why when we discuss about matrix representation of linear maps.
\end{remarkenv}

\begin{theoremenv}
	Let $\{\mathbf{v}_1, \mathbf{v}_2, \dots, \mathbf{v}_n\}$ be a set of elements of a vector space $V$ over a field $\mathcal{K}$, and let $r$ be a positive integer with $r \leq n$.
	We say that $\{\mathbf{v}_1, \mathbf{v}_2, \dots, \mathbf{v}_r\}$ is a \textbf{maximal subset of linearly independent elements} if
	$\mathbf{v}_1, \dots, \mathbf{v}_r$ are linearly independent, and moreover, for any $\mathbf{v}_i$ with $i > r$, the set
	$\{\mathbf{v}_1, \dots, \mathbf{v}_r, \mathbf{v}_i\}$ is linearly dependent.

	If, in addition, $\{\mathbf{v}_1, \mathbf{v}_2, \dots, \mathbf{v}_n\}$ is a set of generators of $V$,
	then the set $\{\mathbf{v}_1, \mathbf{v}_2, \dots, \mathbf{v}_r\}$ is a \textbf{basis} of $V$.
\end{theoremenv}


\section{Dimension}

\begin{theoremenv}
	Let $V$ be a vector space over a field $\mathcal{K}$.
	Let $\{\mathbf{v}_1, \mathbf{v}_2, \dots, \mathbf{v}_m\}$ be a basis of $V$ over $\mathcal{K}$.
	Let $\mathbf{w}_1, \mathbf{w}_2, \dots, \mathbf{w}_n$ be elements of $V$, and assume that $n > m$.
	Then $\mathbf{w}_1, \mathbf{w}_2, \dots, \mathbf{w}_n$ are linearly dependent.
\end{theoremenv}

\begin{theoremenv}\label{thm:dimension_well_defined}
	Let $V$ be a vector space over a field $\mathcal{K}$, and suppose that one basis of $V$ has $n$ elements and another basis has $m$ elements.
	Then $m = n$.
\end{theoremenv}

\begin{remarkenv}
    Theorem~\ref{thm:dimension_well_defined} ensures that the dimension of a finite-dimensional vector space is well-defined, that is, it does not depend on the choice of basis.
\end{remarkenv}

\begin{definitionenv}
	Let $V$ be a vector space over a field $\mathcal{K}$.
	If $V$ has a finite basis $\mathcal{B}$ consisting of $n$ vectors, then $V$ is said to be \textbf{finite-dimensional}, and the number $n$ is called the \textbf{dimension} of $V$, denoted by $\dim(V) = n$.

	If no finite basis exists, $V$ is said to be \textbf{infinite-dimensional}.
\end{definitionenv}

\begin{remarkenv}
	The existence of a basis in a finite-dimensional vector space can be proved via Steinitz Exchange Lemma. However, the existence of a basis in infinite-dimensional vector space relies on the Axiom of Choice.
\end{remarkenv}
\begin{exampleenv}
	Let $\mathcal{K}$ be a field.
	Define
	\[
	V = \mathcal{K}[x] =
	\left\{
	a_0 + a_1x + a_2x^2 + \cdots + a_nx^n
	\ \middle| \
	n \in \mathbb{N},\ a_i \in \mathcal{K}
	\right\},
	\]
	the set of all polynomials with coefficients in $\mathcal{K}$. Then $V$ is a vector space over $\mathcal{K}$ with the usual polynomial addition and scalar multiplication.

	The set
	\[
	\{\, 1, x, x^2, x^3, \dots \}
	\]
	is linearly independent, and no finite subset of it spans $V$.
	Hence $\mathcal{K}[x]$ is an \textbf{infinite-dimensional} vector space over $\mathcal{K}$.
\end{exampleenv}

\begin{definitionenv}
	Let $V$ be a vector space over a field $\mathcal{K}$.
	A subspace of $V$ with dimension $1$ is called a \textbf{line} in $V$.
	A subspace of $V$ with dimension $2$ is called a \textbf{plane} in $V$.
\end{definitionenv}


\begin{definitionenv}
	Let $V$ be a finite-dimensional vector space over a field $\mathcal{K}$, and let
	$\{\mathbf{e}_1, \mathbf{e}_2, \dots, \mathbf{e}_n\}$ be an ordered basis of $V$.
	For any vector $\mathbf{v} \in V$, there exist unique scalars
	$a_1, a_2, \dots, a_n \in \mathcal{K}$ such that
	\[
	\mathbf{v} = a_1 \mathbf{e}_1 + a_2 \mathbf{e}_2 + \cdots + a_n \mathbf{e}_n.
	\]
	The scalars $a_1, a_2, \dots, a_n$ are called the \textbf{components} or \textbf{coordinates} of $\mathbf{v}$
	with respect to the basis $\{\mathbf{e}_1, \mathbf{e}_2, \dots, \mathbf{e}_n\}$. And the $n$-tuple $A=(a_1,\dots,a_n)$ is called the \textbf{coordinate vector} of \textbf{v}.
\end{definitionenv}

\begin{definitionenv}
	Let $V$ be a vector space over a field $\mathcal{K}$, and let
	$S \subseteq V$.
	A subset $\mathcal{B} \subseteq S$ is called a \textbf{maximal linearly independent subset} of $S$ if
	\begin{enumerate}
		\item $\mathcal{B}$ is linearly independent;
		\item for any vector $\mathbf{v} \in S \setminus \mathcal{B}$, the set
		$\mathcal{B} \cup \{\mathbf{v}\}$ is linearly dependent.
	\end{enumerate}
\end{definitionenv}

\begin{theoremenv}
	Let $V$ be a vector space over a field $\mathcal{K}$, and let the set
	$\{\mathbf{v}_1, \mathbf{v}_2, \dots, \mathbf{v}_n\}$ be a maximal set of linearly independent elements of $V$.
	Then $\{\mathbf{v}_1, \mathbf{v}_2, \dots, \mathbf{v}_n\}$ is a basis of $V$.
\end{theoremenv}

% \begin{remarkenv}
% 	You might think, ``Oh! Why does the theorem above look similar to something previous. Is the author mentally unstable?" Actually, they are different because of their starting points.
% \end{remarkenv}

\begin{theoremenv}
	Let $V$ be a vector space over a field $\mathcal{K}$ of dimension $n$, and let
	$\mathbf{v}_1, \mathbf{v}_2, \dots, \mathbf{v}_n$ be linearly independent elements of $V$.
	Then $\mathbf{v}_1, \mathbf{v}_2, \dots, \mathbf{v}_n$ constitute a basis of $V$.
\end{theoremenv}

\begin{propositionenv}
	Let $V$ be a vector space over a field $\mathcal{K}$, and let $W$ be a subspace of $V$.
	If $\dim W = \dim V$, then $V = W$.
\end{propositionenv}

\begin{remarkenv}
	This proposition gives us a way to prove that two sets (vector space) are the same.
\end{remarkenv}

\begin{propositionenv}
	Let $V$ be a vector space over a field $\mathcal{K}$ of dimension $n$.
	Let $r$ be a positive integer with $r < n$, and let
	$\mathbf{v}_1, \mathbf{v}_2, \dots, \mathbf{v}_r$ be linearly independent elements of $V$.
	Then there exist elements $\mathbf{v}_{r+1}, \mathbf{v}_{r+2}, \dots, \mathbf{v}_n$ in $V$ such that
	\[
	\{\mathbf{v}_1, \mathbf{v}_2, \dots, \mathbf{v}_n\}
	\]
	is a basis of $V$.
\end{propositionenv}

\begin{remarkenv}
    These theorems and propositions all point to the same idea: a linearly independent set may be too small to span the whole vector space, while a spanning set may be too large to be linearly independent. A basis lies exactly at the boundary between these two situations. It is both a minimal spanning set and a maximal linearly independent set. In this sense, it plays a role somewhat analogous to a supremum.
\end{remarkenv}

\begin{theoremenv}
	Let $V$ be a vector space over a field $\mathcal{K}$ having a basis consisting of $n$ elements.
	Let $W$ be a subspace of $V$ that does not consist of the zero vector alone.
	Then $W$ has a basis, and the dimension of $W$ satisfies
	\[
	\dim W \leq n.
	\]
\end{theoremenv}

\newpage